% Lecture 4 notes: \input by notes.tex, \subimport'ed by the semester book.
% Exam questions (зачёт): 7 (линейная и циклическая свёртки, линейная через циклическую),
% 8 (теорема о свёртке), 9 (умножение с помощью БПФ, оценка сложности).
\begin{lecture}{4}{Свёртки и умножение с помощью БПФ}

\section{Зачем нужны свёртки}

В лекции~2 мы видели, что произведение чисел $a = \sum_i a_i B^i$ и $b = \sum_j b_j B^j$
выражается через суммы $\sum_{i+j=t} a_i b_j$ "--- \emph{свёртку} последовательностей цифр.
Школьный метод вычисляет свёртку «в лоб» за $\Theta(N^2)$ операций, методы Карацубы и
Тоома"--~Кука "--- за $\bigO(N^{1{,}585})$ и $\bigO(N^{1+\varepsilon})$. В этой лекции мы
построим алгоритм, который вычисляет свёртку длины $N$ за $\bigO(N \log N)$ операций в подходящем
кольце: \emph{быстрое преобразование Фурье} (БПФ) переводит свёртку в поточечное произведение.

Свёртки встречаются в криптографии не только через умножение длинных чисел (RSA,
Диффи"--~Хеллман). Постквантовые схемы на решётках, стандартизованные NIST в 2024~г.
(ML-KEM/Kyber, ML-DSA/Dilithium), работают в кольце многочленов $\Zn{q}[t]/(t^{256} + 1)$:
умножение в нём "--- это \emph{негациклическая} свёртка длины $256$, и вычисляют её именно
быстрым преобразованием над конечным полем (теоретико-числовым преобразованием, лекция~5).

\section{Свёртки}

Всюду в этом разделе $R$ "--- коммутативное кольцо с единицей (нам будут нужны $\ZZ$, $\CC$
и $\GF{p}$), последовательности $x = (x_0, \dots, x_{M-1})$ и $y = (y_0, \dots, y_{N-1})$
имеют элементы из $R$. Последовательности удобно отождествлять с многочленами:
\[
  x(t) = \sum_{i=0}^{M-1} x_i t^i, \qquad y(t) = \sum_{j=0}^{N-1} y_j t^j.
\]

\begin{definition}\label{l4:linear}
\emph{Линейной} (ациклической) \emph{свёрткой} последовательностей $x$ и $y$ называется
последовательность $z = x * y$ длины $M + N - 1$,
\[
  z_k = \sum_{i + j = k} x_i y_j = \sum_{i = \max(0,\, k - N + 1)}^{\min(k,\, M - 1)} x_i\, y_{k-i},
  \qquad k = 0, 1, \dots, M + N - 2.
\]
\end{definition}

Раскрывая скобки в $x(t)\,y(t)$ и приводя подобные, получаем коэффициент при $t^k$, равный
$\sum_{i+j=k} x_i y_j$. Итак, \emph{линейная свёртка "--- это последовательность коэффициентов
произведения многочленов}: $(x * y)(t) = x(t)\,y(t)$. Отсюда, в частности, свёртка коммутативна,
ассоциативна и билинейна. Вычисление по определению требует $MN$ умножений.

\begin{proposition}\label{l4:intmul}
Пусть $a = \sum_{i=0}^{M-1} x_i B^i$ и $b = \sum_{j=0}^{N-1} y_j B^j$ "--- записи чисел по
основанию $B$, $z = x * y$. Тогда
\[
  ab = \sum_{k=0}^{M+N-2} z_k B^k, \qquad 0 \le z_k \le \min(M, N)\,(B - 1)^2 .
\]
Цифры произведения получаются из $z_k$ переносами за $\bigO(M + N)$ операций.
\end{proposition}

\begin{proof}
Подстановка $t = B$ "--- гомоморфизм колец $\ZZ[t] \to \ZZ$, поэтому
$ab = x(B)\,y(B) = z(B)$. В сумме $z_k = \sum_{i+j=k} x_i y_j$ не более $\min(M, N)$ слагаемых
(индекс $i$ однозначно определяет $j$), и каждое не больше $(B-1)^2$.

Переносы: $c_0 = 0$, $t_k = z_k + c_k$, цифра $w_k = t_k \bmod B$, $c_{k+1} = \floor{t_k / B}$.
Обозначим $K = \min(M, N)$; индукцией по $k$ проверим, что $c_k < KB$: если это так, то
$t_k < K(B-1)^2 + KB = K(B^2 - B + 1) \le KB^2$, откуда $c_{k+1} < KB$. Значит, переносы и все
промежуточные величины занимают $\bigO(1)$ машинных слов (при $K$, ограниченном размером слова),
и проход по $M + N - 1$ коэффициентам (плюс $\bigO(1)$ дополнительных разрядов, так как
у $ab$ не более $M + N$ цифр) стоит $\bigO(M + N)$ операций.
\end{proof}

\begin{example}\label{l4:conv-ex}
$a = 572$, $b = 38$; цифры (младшие первыми) $x = (2, 7, 5)$, $y = (8, 3)$:
\[
  z_0 = 2 \cdot 8 = 16, \quad z_1 = 2 \cdot 3 + 7 \cdot 8 = 62, \quad
  z_2 = 7 \cdot 3 + 5 \cdot 8 = 61, \quad z_3 = 5 \cdot 3 = 15.
\]
Переносы: $16 \to$ цифра $6$, перенос $1$; $62 + 1 = 63 \to 3$, перенос $6$; $61 + 6 = 67 \to 7$,
перенос $6$; $15 + 6 = 21 \to 1$, перенос $2$; итак, $572 \cdot 38 = 21736$.
\end{example}

\subsection{Циклическая и негациклическая свёртки}

Преобразование Фурье естественно работает не с линейной, а с \emph{циклической} свёрткой:
индексы складываются по модулю длины.

\begin{definition}\label{l4:cyclic}
Пусть $x = (x_0, \dots, x_{n-1})$, $y = (y_0, \dots, y_{n-1})$. \emph{Циклической свёрткой}
длины $n$ называется последовательность $x \circledast y$,
\[
  (x \circledast y)_k = \sum_{i + j \equiv k \ (\mathrm{mod}\ n)} x_i y_j
  = \sum_{i=0}^{n-1} x_i\, y_{(k - i) \bmod n}, \qquad k = 0, \dots, n - 1,
\]
а \emph{негациклической} "--- последовательность
\[
  (x \circledast_{-} y)_k = \sum_{i + j = k} x_i y_j - \sum_{i + j = k + n} x_i y_j .
\]
\end{definition}

\begin{lemma}\label{l4:fold}
Пусть $x$, $y$ "--- последовательности длины $n$ и $z = x * y$ (длины $2n - 1$; положим
$z_{2n-1} = 0$). Тогда
\[
  (x \circledast y)_k = z_k + z_{k+n}, \qquad (x \circledast_{-} y)_k = z_k - z_{k+n}, \qquad
  0 \le k < n.
\]
Иначе говоря, $(x \circledast y)(t) \equiv x(t)\,y(t) \pmod{t^n - 1}$ и
$(x \circledast_{-} y)(t) \equiv x(t)\,y(t) \pmod{t^n + 1}$, где слева стоят многочлены степени
меньше $n$.
\end{lemma}

\begin{proof}
При $0 \le i, j < n$ сумма $i + j$ лежит в $[0, 2n - 2]$, поэтому $i + j \equiv k \pmod n$
равносильно $i + j \in \{k, k + n\}$. Разбивая сумму из определения на две части, получаем
$z_k + z_{k+n}$ (и $z_k - z_{k+n}$ для негациклической). Для многочленов:
$x(t)y(t) = \sum_{k<n} z_k t^k + t^n \sum_{k<n} z_{k+n} t^k$, и $t^n \equiv \pm 1$ по модулю
$t^n \mp 1$.
\end{proof}

Таким образом, циклическая свёртка «заворачивает» старшие коэффициенты линейной на младшие
позиции. Если заворачиваться нечему, свёртки совпадают "--- это главный приём, сводящий
линейную свёртку к циклической.

\begin{theorem}[дополнение нулями]\label{l4:zeropad}
Пусть $x$ и $y$ имеют длины $M$ и $N$, $L \ge M + N - 1$, и $\bar x$, $\bar y$ "--- последовательности
длины $L$, полученные из $x$, $y$ дописыванием нулей. Тогда
\[
  (\bar x \circledast \bar y)_k = (x * y)_k \quad (0 \le k \le M + N - 2), \qquad
  (\bar x \circledast \bar y)_k = 0 \quad (M + N - 1 \le k < L).
\]
Если же $L < M + N - 1$, то по циклической свёртке длины $L$ линейную, вообще говоря,
не восстановить.
\end{theorem}

\begin{proof}
Слагаемое $\bar x_i \bar y_j$ может быть ненулевым только при $i < M$, $j < N$, т.\,е. при
$i + j \le M + N - 2 < L$. Для таких пар условие $i + j \equiv k \pmod L$ с $0 \le k < L$
равносильно $i + j = k$, так что ненулевые слагаемые в $(\bar x \circledast \bar y)_k$ и
в $(x * y)_k$ одни и те же. (Это частный случай леммы~\ref{l4:fold}: все $z_{k+L} = 0$.)

Если $L < M + N - 1$, то по лемме~\ref{l4:fold} старший коэффициент
$z_{M+N-2} = x_{M-1} y_{N-1}$ складывается с $z_{M+N-2-L}$, и разные линейные свёртки могут дать
одну и ту же циклическую. Например, при $M = N = L = 2$ над $\ZZ$ $(1, 1) \circledast (1, 1) = (2, 2) = (2, 0) \circledast (1, 1)$, хотя
$(1, 1) * (1, 1) = (1, 2, 1) \ne (2, 2, 0) = (2, 0) * (1, 1)$.
\end{proof}

\begin{example}
Для $x = (2, 7, 5)$, $y = (8, 3)$ из примера~\ref{l4:conv-ex} линейная свёртка
$z = (16, 62, 61, 15)$, $M + N - 1 = 4$. Циклическая свёртка длины $4$ (с нулями
$\bar x = (2, 7, 5, 0)$, $\bar y = (8, 3, 0, 0)$) совпадает с $z$. Длины $3$ уже не хватает:
\[
  x \circledast y = (16 + 15,\ 62,\ 61) = (31, 62, 61), \qquad
  x \circledast_{-} y = (16 - 15,\ 62,\ 61) = (1, 62, 61).
\]
Это согласуется с многочленами: $572 \cdot 38 = 21736$, а
$31 + 62 \cdot 10 + 61 \cdot 100 = 6751 = 21736 \bmod 999$ (значение при $t = 10$ по модулю
$t^3 - 1 = 999$) и $1 + 620 + 6100 = 6721 = 21736 \bmod 1001$ (по модулю $t^3 + 1 = 1001$).
\end{example}

\taskexample{conv-wrap}

\checkyourself{Вычислите циклическую и негациклическую свёртки длины $3$ последовательностей
$(1, 2, 3)$ и $(4, 5, 6)$. Как, зная обе, найти линейную свёртку?}

\checkyourself{Почему для сведения линейной свёртки к циклической длины $L$ нужно именно
$L \ge M + N - 1$, а не $L \ge \max(M, N)$?}

\section{Дискретное преобразование Фурье}

\subsection{Главные корни из единицы}

\begin{definition}\label{l4:principal}
Элемент $\omega$ кольца $R$ называется \emph{главным корнем степени $n$ из единицы}, если
\[
  \omega^n = 1 \qquad\text{и}\qquad \sum_{j=0}^{n-1} \omega^{jk} = 0 \quad\text{при } 0 < k < n.
\]
\end{definition}

Второе условие "--- \emph{условие ортогональности}; на нём держится обратимость преобразования.
Заметим, что главный корень имеет порядок ровно $n$: если $\omega^k = 1$ при $0 < k < n$, то
$\sum_j \omega^{jk} = n \cdot 1$, а это не ноль, если $n$ обратимо в $R$.

\begin{proposition}\label{l4:field}
Пусть $R$ "--- поле и $n \cdot 1 \ne 0$ в $R$. Тогда $\omega$ "--- главный корень
степени $n$ из единицы тогда и только тогда, когда порядок $\omega$ в мультипликативной группе
равен $n$. В частности:
\begin{enumerate}
  \item в $\CC$ главные корни степени $n$ "--- это первообразные корни $e^{2\pi i s/n}$,
    $\gcd(s, n) = 1$, например $\omega = e^{2\pi i/n}$;
  \item в $\GF{p}$ главный корень степени $n$ существует тогда и только тогда, когда
    $n \mid p - 1$; если $g$ "--- первообразный корень по модулю $p$, то
    $\omega = g^{(p-1)/n}$ "--- главный корень степени $n$.
\end{enumerate}
\end{proposition}

\begin{proof}
Пусть порядок $\omega$ равен $n$ и $0 < k < n$. Тогда $\omega^k \ne 1$, и по формуле суммы
геометрической прогрессии
\[
  (\omega^k - 1) \sum_{j=0}^{n-1} \omega^{jk} = \omega^{kn} - 1 = 0,
\]
а так как $\omega^k - 1 \ne 0$ и в поле нет делителей нуля, сумма равна нулю. Обратно, если
$\omega^n = 1$, но порядок $d$ элемента $\omega$ меньше $n$, то при $k = d$ все слагаемые
$\omega^{jd}$ равны $1$ и сумма равна $n \cdot 1 \ne 0$, т.\,е. $\omega$ не главный корень.

Пункт~1: группа корней степени $n$ в $\CC$ циклическая, порождена $e^{2\pi i/n}$.
Пункт~2: при $n \mid p - 1$ имеем $n < p$, так что $n \cdot 1 \ne 0$ в $\GF{p}$; группа $\Znx{p}$ циклическая порядка $p - 1$; элемент порядка $n$ в ней есть тогда и
только тогда, когда $n \mid p - 1$, и $g^{(p-1)/n}$ имеет порядок ровно $n$.
\end{proof}

В кольцах с делителями нуля одного порядка $n$ мало. Например, в $\Zn{15}$ элемент $4$ имеет
порядок $2$ ($4^2 = 16 \equiv 1$), но $1 + 4 = 5 \not\equiv 0$, так что $4$ не является главным
корнем степени $2$. Нас будут интересовать длины $n = 2^m$, и для них есть удобный критерий.

\begin{lemma}\label{l4:pow2}
Пусть $n = 2^m$, $m \ge 1$, и $2$ обратимо в $R$. Элемент $\omega$ "--- главный корень степени
$n$ из единицы тогда и только тогда, когда $\omega^{n/2} = -1$.
\end{lemma}

\begin{proof}
Многочлены $\sum_{j=0}^{2^m - 1} s^j$ и $\prod_{i=0}^{m-1} (1 + s^{2^i})$ совпадают: раскрывая
скобки в произведении, получаем $s^j$ по одному разу для каждого $j < 2^m$ (двоичная запись
$j$ указывает, из каких скобок взят множитель $s^{2^i}$).

($\Leftarrow$) Пусть $\omega^{n/2} = -1$; тогда $\omega^n = 1$. Возьмём $0 < k < n$ и запишем
$k = 2^r u$, $u$ нечётно, $r < m$. При $i = m - 1 - r$ получаем
$\omega^{k \cdot 2^i} = \omega^{u \cdot 2^{m-1}} = (\omega^{n/2})^u = -1$, поэтому множитель
$1 + \omega^{k 2^i}$ равен нулю, и $\sum_j \omega^{jk} = \prod_i (1 + \omega^{k 2^i}) = 0$.

($\Rightarrow$) Положим $\eta = \omega^{n/2}$, $\eta^2 = 1$. Условие ортогональности при
$k = n/2$ даёт $0 = \sum_{j=0}^{n-1} \eta^j = \frac n2 (1 + \eta)$. Так как $2$ обратимо,
обратимо и $\frac n2 = 2^{m-1}$, откуда $\eta = -1$.
\end{proof}

Например, в $\GF{17}$: $4^2 = 16 = -1$, так что $\omega = 4$ "--- главный корень степени $4$;
а $\omega = 16 = -1$ удовлетворяет $\omega^4 = 1$, но $\omega^2 = 1 \ne -1$ "--- это не главный
корень степени $4$. В $\GF{257}$, $257 = 2^8 + 1$: $16^2 = 256 = -1$, $4^4 = -1$, $2^8 = -1$,
т.\,е. $16$, $4$ и $2$ "--- главные корни степеней $4$, $8$ и $16$ соответственно.

\subsection{Прямое и обратное преобразования}

\begin{definition}\label{l4:dft}
Пусть $\omega \in R$ "--- главный корень степени $n$ из единицы. \emph{Дискретным
преобразованием Фурье} (ДПФ) последовательности $x = (x_0, \dots, x_{n-1})$ называется
последовательность $X = \mathrm{DFT}_\omega(x)$,
\[
  X_k = \sum_{j=0}^{n-1} x_j\, \omega^{jk} = x(\omega^k), \qquad k = 0, \dots, n - 1.
\]
\end{definition}

ДПФ "--- это вычисление многочлена $x(t)$ в $n$ точках $1, \omega, \dots, \omega^{n-1}$,
т.\,е. умножение вектора на матрицу Вандермонда $W = (\omega^{jk})_{j,k}$. Над $\CC$ часто
берут $\omega = e^{-2\pi i/n}$; для теории знак показателя не важен. Вычисление по
определению требует $\Theta(n^2)$ операций в $R$.

\begin{theorem}[обратное ДПФ]\label{l4:inverse}
Пусть $\omega$ "--- главный корень степени $n$ из единицы и $n$ обратимо в $R$. Тогда
$\omega^{-1}$ тоже главный корень степени $n$, и для $X = \mathrm{DFT}_\omega(x)$
\[
  x_j = \frac1n \sum_{k=0}^{n-1} X_k\, \omega^{-jk}, \qquad\text{т.\,е.}\qquad
  x = n^{-1}\, \mathrm{DFT}_{\omega^{-1}}(X).
\]
\end{theorem}

\begin{proof}
$\omega$ обратим, так как $\omega \cdot \omega^{n-1} = 1$, и $\omega^{-1} = \omega^{n-1}$.
Для $0 < k < n$: $\sum_j \omega^{-jk} = \sum_j \omega^{j(n-k)} = 0$, так как $0 < n - k < n$;
значит, $\omega^{-1}$ "--- главный корень.

Сначала докажем \emph{соотношение ортогональности}: для целого $s$
\[
  \sum_{k=0}^{n-1} \omega^{sk} = \begin{cases} n, & n \mid s,\\ 0, & n \nmid s. \end{cases}
\]
Действительно, $\omega^{sk} = \omega^{rk}$, где $r = s \bmod n$ (так как $\omega^n = 1$); при $r = 0$
все слагаемые равны $1$, а при $0 < r < n$ сумма равна нулю по определению главного корня.
Теперь
\[
  \sum_{k=0}^{n-1} X_k \omega^{-jk} = \sum_{k} \sum_{i} x_i \omega^{ik} \omega^{-jk}
  = \sum_{i=0}^{n-1} x_i \sum_{k=0}^{n-1} \omega^{(i - j)k} = n\, x_j,
\]
так как при $0 \le i, j < n$ число $i - j$ делится на $n$ только при $i = j$.
\end{proof}

Итак, обратное ДПФ "--- это то же ДПФ с корнем $\omega^{-1}$ и умножением на $n^{-1}$; любой
алгоритм быстрого вычисления ДПФ годится и для обратного. На языке многочленов теорема
утверждает, что многочлен степени меньше $n$ восстанавливается по значениям в $n$ точках
$\omega^k$ "--- это интерполяция, как в методе Тоома"--~Кука (лекция~2), но с особенно удобными
точками.

\checkyourself{Где в доказательстве теоремы~\ref{l4:inverse} используется условие
ортогональности, а где "--- обратимость $n$? Приведите пример кольца, где $\omega^n = 1$,
порядок $\omega$ равен $n$, но обратного преобразования нет.}

\section{Теорема о свёртке}

\begin{theorem}[о свёртке]\label{l4:convthm}
Пусть $\omega \in R$, $\omega^n = 1$, и $x$, $y$ "--- последовательности длины $n$. Тогда
\[
  \mathrm{DFT}_\omega(x \circledast y) = \mathrm{DFT}_\omega(x) \cdot \mathrm{DFT}_\omega(y)
\]
(произведение поточечное: $Z_k = X_k Y_k$). Если, кроме того, $\omega$ "--- главный корень и $n$
обратимо в $R$, то
\[
  x \circledast y = n^{-1}\, \mathrm{DFT}_{\omega^{-1}}\bigl(\mathrm{DFT}_\omega(x) \cdot
  \mathrm{DFT}_\omega(y)\bigr).
\]
\end{theorem}

\begin{proof}
Пусть $z = x \circledast y$. Поскольку $\omega^n = 1$, степень $\omega^{lk}$ зависит лишь от
$l \bmod n$, и
\[
  Z_k = \sum_{l=0}^{n-1} \Bigl(\sum_{i + j \equiv l\ (\mathrm{mod}\ n)} x_i y_j\Bigr) \omega^{lk}
  = \sum_{i=0}^{n-1} \sum_{j=0}^{n-1} x_i y_j\, \omega^{(i + j)k}
  = \Bigl(\sum_i x_i \omega^{ik}\Bigr) \Bigl(\sum_j y_j \omega^{jk}\Bigr) = X_k Y_k .
\]
(Каждая пара $(i, j)$ входит ровно в одну внутреннюю сумму "--- с $l = (i + j) \bmod n$.)
Вторая формула получается применением теоремы~\ref{l4:inverse} к $z$.
\end{proof}

\begin{remark}
То же доказательство на языке многочленов: по лемме~\ref{l4:fold}
$z(t) = x(t)y(t) - q(t)(t^n - 1)$, а $\omega^k$ "--- корень $t^n - 1$, поэтому
$z(\omega^k) = x(\omega^k)\,y(\omega^k)$. Значения произведения в точках "--- произведения
значений; в этом и состоит идея: \emph{вычислить в точках "--- перемножить "--- интерполировать}.
\end{remark}

\begin{corollary}\label{l4:linconv}
Пусть $L \ge M + N - 1$, $\omega$ "--- главный корень степени $L$ и $L$ обратимо в $R$. Тогда
линейная свёртка последовательностей длин $M$ и $N$ равна первым $M + N - 1$ элементам
$L^{-1} \mathrm{DFT}_{\omega^{-1}}(\mathrm{DFT}_\omega(\bar x) \cdot \mathrm{DFT}_\omega(\bar y))$,
где $\bar x$, $\bar y$ "--- дополненные нулями до длины $L$ последовательности.
\end{corollary}

\begin{proof}
Теорема~\ref{l4:zeropad} и теорема~\ref{l4:convthm}.
\end{proof}

\begin{example}\label{l4:dft-c}
Над $\CC$ при $n = 4$ возьмём $\omega = i$; матрица ДПФ
$W = \left(\begin{smallmatrix} 1 & 1 & 1 & 1\\ 1 & i & -1 & -i\\ 1 & -1 & 1 & -1\\ 1 & -i & -1 & i
\end{smallmatrix}\right)$. Для $\bar x = (2, 7, 5, 0)$, $\bar y = (8, 3, 0, 0)$:
\[
  X = (14,\ -3 + 7i,\ 0,\ -3 - 7i), \qquad Y = (11,\ 8 + 3i,\ 5,\ 8 - 3i),
\]
\[
  Z = XY = (154,\ -45 + 47i,\ 0,\ -45 - 47i).
\]
Обратное ДПФ: $z_0 = \frac14 (154 + (-45 + 47i) + 0 + (-45 - 47i)) = \frac{64}{4} = 16$,
$z_1 = \frac14\bigl(154 + (-45 + 47i)(-i) + (-45 - 47i)\,i\bigr) = \frac{154 + 94}{4} = 62$,
аналогично $z_2 = 61$, $z_3 = 15$ "--- линейная свёртка из примера~\ref{l4:conv-ex}.
\end{example}

\begin{proposition}[негациклическая свёртка через циклическую]\label{l4:dwt}
Пусть $\psi \in R$, $\psi^n = -1$, $x$, $y$ "--- последовательности длины $n$. Положим
$x'_j = \psi^j x_j$, $y'_j = \psi^j y_j$. Тогда $(x' \circledast y')_k = \psi^k (x \circledast_{-} y)_k$.
\end{proposition}

\begin{proof}
$(x' \circledast y')_k = \sum_{i+j=k} \psi^{k} x_i y_j + \sum_{i+j=k+n} \psi^{k+n} x_i y_j
= \psi^k \bigl(\sum_{i+j=k} x_i y_j - \sum_{i+j=k+n} x_i y_j\bigr)$, так как $\psi^n = -1$.
\end{proof}

Если $\psi$ "--- главный корень степени $2n$, то $\psi^n = -1$ (лемма~\ref{l4:pow2}) и
$\omega = \psi^2$ "--- главный корень степени $n$; негациклическая свёртка вычисляется тремя ДПФ
длины $n$ и $\bigO(n)$ умножениями на $\psi^{\pm j}$ (\emph{взвешенное} преобразование,
\cite[разд.~9.5.4]{crandall-pomerance}). Так умножают в кольцах $\Zn{q}[t]/(t^n + 1)$ (Kyber) и
по модулю чисел Ферма $2^n + 1$ (лекция~5).

\taskexample{dft-f17}

\checkyourself{Теорема о свёртке (первая часть) использует только $\omega^n = 1$. Почему же
в $\GF{17}$ нельзя умножать с $\omega = 16$ при $n = 4$?}

\section{Быстрое преобразование Фурье}

\subsection{Алгоритм Кули"--~Тьюки}

Пусть $n = 2m$ чётно и $\omega$ "--- главный корень степени $n$ с $\omega^{m} = -1$ (для
$n = 2^r$ это выполнено по лемме~\ref{l4:pow2}). Разобьём сумму в $X_k$ на слагаемые с чётными
и нечётными индексами:
\[
  X_k = \sum_{j=0}^{m-1} x_{2j}\, (\omega^2)^{jk} + \omega^k \sum_{j=0}^{m-1} x_{2j+1}\, (\omega^2)^{jk}
  = E_k + \omega^k O_k,
\]
где $E = \mathrm{DFT}_{\omega^2}(x_0, x_2, \dots, x_{n-2})$ и
$O = \mathrm{DFT}_{\omega^2}(x_1, x_3, \dots, x_{n-1})$ "--- ДПФ \emph{половинной} длины $m$
(элемент $\omega^2$ имеет порядок $m$, и последовательности $E_k$, $O_k$ $m$-периодичны по $k$).
Поскольку $\omega^{k + m} = -\omega^k$,
\[
  \boxed{\;X_k = E_k + \omega^k O_k, \qquad X_{k+m} = E_k - \omega^k O_k, \qquad 0 \le k < m.\;}
\]
Пара формул называется \emph{бабочкой} (рис.~\ref{l4:fig-butterfly}): по двум входам $E_k$, $O_k$
одно умножение на $\omega^k$ («поворачивающий множитель») и одно сложение с вычитанием дают
два выхода. В~таком виде алгоритм опубликовали Дж.~Кули и Дж.~Тьюки в 1965~г.~\cite{cooley-tukey}
(идея восходит к Гауссу).

\begin{figure}[ht]
\centering
\begin{tikzpicture}[x=1cm, y=1cm, >={Stealth[length=2mm]}, font=\small]
  % single butterfly
  \node[left] (E) at (0, 1) {$E_k$};
  \node[left] (O) at (0, 0) {$O_k$};
  \node[right] (X1) at (3, 1) {$E_k + \omega^k O_k = X_k$};
  \node[right] (X2) at (3, 0) {$E_k - \omega^k O_k = X_{k+m}$};
  \draw[->] (0, 1) -- (3, 1);
  \draw[->] (0, 0) -- node[pos=0.18, below] {$\omega^k$} (3, 1);
  \draw[->] (0, 1) -- (3, 0);
  \draw[->] (0, 0) -- node[pos=0.85, below] {$-1$} (3, 0);
  % 4-point FFT flow graph
  \begin{scope}[xshift=8.6cm, yshift=-0.9cm, y=0.9cm]
    \foreach \lab [count=\r from 0] in {x_0, x_2, x_1, x_3} {
      \node[left] at (0, 3 - \r) {$\lab$};
    }
    \foreach \lab [count=\r from 0] in {X_0, X_1, X_2, X_3} {
      \node[right] at (4, 3 - \r) {$\lab$};
    }
    \node[above] at (2, 3) {$E_0$};
    \node[above, fill=white, inner sep=1pt] at (2, 2.12) {$E_1$};
    \node[above, fill=white, inner sep=1pt] at (2, 1.12) {$O_0$};
    \node[below] at (2, 0) {$O_1$};
    % stage 1: (x0,x2) -> (E0,E1), (x1,x3) -> (O0,O1)
    \foreach \a/\b in {3/2, 1/0} {
      \draw[->] (0, \a) -- (2, \a); \draw[->] (0, \b) -- (2, \a);
      \draw[->] (0, \a) -- (2, \b); \draw[->] (0, \b) -- (2, \b);
    }
    % stage 2: (E0,O0) -> (X0,X2), (E1,O1) -> (X1,X3)
    \draw[->] (2, 3) -- (4, 3); \draw[->] (2, 1) -- (4, 3);
    \draw[->] (2, 3) -- (4, 1); \draw[->] (2, 1) -- (4, 1);
    \draw[->] (2, 2) -- (4, 2); \draw[->] (2, 0) -- (4, 2);
    \draw[->] (2, 2) -- (4, 0); \draw[->] (2, 0) -- (4, 0);
    \node[below right, inner sep=1pt] at (2.05, 0.05) {$\omega$};
    \foreach \y in {0, 1, 2, 3} { \fill (2, \y) circle (1.2pt); }
  \end{scope}
\end{tikzpicture}
\caption{Бабочка (слева) и граф БПФ длины $4$ (справа); на нижних рёбрах каждой бабочки "---
множитель $-1$, на рёбрах от $O_k$ "--- множитель $\omega^k$.}
\label{l4:fig-butterfly}
\end{figure}

\begin{algorithm}[БПФ, прореживание по времени]\label{l4:alg-fft}
\begin{algorithmic}[1]
\Require $n = 2^r$, $x = (x_0, \dots, x_{n-1}) \in R^n$, главный корень $\omega$ степени $n$.
\Ensure $X = \mathrm{DFT}_\omega(x)$.
\Function{FFT}{$x, \omega$}
  \If{$n = 1$} \Return $(x_0)$ \EndIf
  \State $E \gets \Call{FFT}{(x_0, x_2, \dots, x_{n-2}),\ \omega^2}$
  \State $O \gets \Call{FFT}{(x_1, x_3, \dots, x_{n-1}),\ \omega^2}$
  \State $\tau \gets 1$
  \For{$k = 0, 1, \dots, n/2 - 1$} \Comment{$\tau = \omega^k$}
    \State $u \gets \tau \cdot O_k$
    \State $X_k \gets E_k + u$, \quad $X_{k + n/2} \gets E_k - u$
    \State $\tau \gets \tau \cdot \omega$
  \EndFor
  \State \Return $X$
\EndFunction
\end{algorithmic}
\end{algorithm}

\begin{theorem}\label{l4:fft}
Пусть $n = 2^r$, $2$ обратимо в $R$, $\omega$ "--- главный корень степени $n$. Алгоритм~\ref{l4:alg-fft}
вычисляет $\mathrm{DFT}_\omega(x)$, выполняя $\frac n2 \log_2 n$ умножений на степени $\omega$ и
$n \log_2 n$ сложений и вычитаний в $R$ (плюс $\bigO(n)$ умножений на вычисление степеней
$\omega$ на каждом уровне, т.\,е. $\bigO(n \log n)$ операций в $R$ всего).
\end{theorem}

\begin{proof}
\emph{Корректность} "--- индукция по $r$. При $r = 0$ ДПФ длины $1$ тождественно. При $r \ge 1$
по лемме~\ref{l4:pow2} $\omega^{n/2} = -1$; если $n \ge 4$, то $(\omega^2)^{n/4} = -1$, и по той же
лемме $\omega^2$ "--- главный корень степени $n/2$ (при $n = 2$ имеем $\omega^2 = 1$, и ДПФ
длины $1$ тривиально), так что рекурсивные вызовы корректны по предположению индукции.
Формулы бабочки выведены выше.

\emph{Сложность}. Пусть $A(n)$ и $S(n)$ "--- числа умножений на $\tau$ и сложений/вычитаний.
Тогда $A(1) = S(1) = 0$ и
\[
  A(n) = 2A(n/2) + \tfrac n2, \qquad S(n) = 2S(n/2) + n .
\]
Индукцией: $A(n) = 2 \cdot \frac n4 (r - 1) + \frac n2 = \frac n2 r$, $S(n) = 2 \cdot \frac n2 (r - 1) + n = nr$.
Обновления $\tau \gets \tau\omega$ дают ещё $\frac n2$ умножений на каждом из $r$ уровней (их можно
заменить одной таблицей степеней $\omega^k$, $k < n$, вычисляемой за $n$ умножений). По основной
теореме (лекция~2, случай~2) то же получаем из $T(n) = 2T(n/2) + \Theta(n) = \Theta(n \log n)$.
\end{proof}

Для сравнения, при $n = 2^{20}$ вычисление ДПФ по определению требует около $10^{12}$
умножений, а БПФ "--- около $10^{7}$.

\begin{example}\label{l4:fft-ex}
В $\GF{257}$ возьмём $n = 4$, $\omega = 16$ ($16^2 = 256 = -1$); $\omega^{-1} = 241 = -16$,
$4^{-1} = 193$. Для $\bar x = (2, 7, 5, 0)$: $E = \mathrm{DFT}_{-1}(x_0, x_2) = (2 + 5, 2 - 5) = (7, 254)$,
$O = (x_1 + x_3, x_1 - x_3) = (7, 7)$, $\omega O_1 = 112$, и бабочки дают
\[
  X = (7 + 7,\ 254 + 112,\ 7 - 7,\ 254 - 112) = (14, 109, 0, 142).
\]
Аналогично для $\bar y = (8, 3, 0, 0)$: $E = (8, 8)$, $O = (3, 3)$, $\omega O_1 = 48$,
$Y = (11, 56, 5, 217)$. Поточечно: $Z = (154, 193, 0, 231)$. Обратное БПФ с $\omega^{-1} = 241$:
$E = (154 + 0, 154 - 0) = (154, 154)$, $O = (193 + 231, 193 - 231) = (167, 219)$,
$\omega^{-1} O_1 = 241 \cdot 219 \equiv 94$, откуда
\[
\begin{gathered}
  u = (154 + 167,\ 154 + 94,\ 154 - 167,\ 154 - 94) \equiv (64, 248, 244, 60),\\
  z = u / 4 = (16, 62, 61, 15).
\end{gathered}
\]
Это снова линейная свёртка из примера~\ref{l4:conv-ex}: все её коэффициенты меньше $257$,
поэтому вычисление в $\GF{257}$ даёт их точно. Сравните с примером~\ref{l4:dft-c}: элемент $16$
играет в $\GF{257}$ роль $i$ ($16^2 = -1$), и $X_1 = -3 + 7 \cdot 16 = 109$.
\end{example}

\subsection{Нерекурсивная реализация}

Раскрывая рекурсию, видим, что на нижнем уровне элементы $x_j$ стоят в порядке
\emph{двоично-обращённых} индексов: при $n = 8$ это $x_0, x_4, x_2, x_6, x_1, x_5, x_3, x_7$
(индекс $j = (j_2 j_1 j_0)_2$ стоит на месте $(j_0 j_1 j_2)_2$). Поэтому БПФ обычно реализуют
без рекурсии и без дополнительной памяти: переставляют элементы массива в двоично-обращённом
порядке, а затем выполняют $\log_2 n$ проходов бабочек с шагами $1, 2, 4, \dots, n/2$
(\cite[гл.~30]{clrs}, \cite[разд.~9.5.2]{crandall-pomerance}).

\checkyourself{Сколько умножений на степени $\omega$ выполняет БПФ длины $1024$? Во сколько
раз это меньше, чем $1024^2$?}

\checkyourself{В каком порядке стоят входы БПФ длины $16$ на нижнем уровне рекурсии?}

\section{Умножение целых чисел с помощью БПФ}

\subsection{Алгоритм}

Соберём всё вместе. По утверждению~\ref{l4:intmul} произведение $ab$ получается переносами из
линейной свёртки цифр, по теореме~\ref{l4:zeropad} линейная свёртка "--- это циклическая свёртка
дополненных нулями последовательностей, а по теореме~\ref{l4:convthm} циклическая свёртка "---
это три ДПФ и поточечное произведение.

\begin{algorithm}[умножение с помощью БПФ]\label{l4:alg-fftmul}
\begin{algorithmic}[1]
\Require $a = \sum_{i<M} x_i B^i$, $b = \sum_{j<N} y_j B^j$; кольцо $R$ ($\CC$ или $\GF{p}$) с главным
  корнем $\omega$ степени $L = 2^r \ge M + N - 1$.
\Ensure $ab$.
\State $\bar x \gets (x_0, \dots, x_{M-1}, 0, \dots, 0)$, $\bar y \gets (y_0, \dots, y_{N-1}, 0, \dots, 0)$
  \Comment{длина $L$}
\State $X \gets \Call{FFT}{\bar x, \omega}$, \quad $Y \gets \Call{FFT}{\bar y, \omega}$
\State $Z_k \gets X_k Y_k$, \ $k = 0, \dots, L - 1$
\State $z \gets L^{-1} \cdot \Call{FFT}{Z, \omega^{-1}}$
  \Comment{над $\CC$: $z_k \gets$ ближайшее целое к $\mathrm{Re}\, z_k$}
\State \Return $\sum_k z_k B^k$, вычисленная переносами
\end{algorithmic}
\end{algorithm}

Коэффициенты $z_k$ "--- \emph{целые} числа из отрезка $[0, \min(M, N)(B-1)^2]$, а алгоритм
вычисляет их образы в $R$. Чтобы восстановить их однозначно, нужно одно из двух.

\begin{itemize}
  \item \textbf{Точная арифметика в $\GF{p}$.} Если $p > \min(M, N)(B - 1)^2$, то $z_k$ "---
    это просто наименьший неотрицательный вычет. Кроме того, нужно $L \mid p - 1$
    (предложение~\ref{l4:field}). Удобны простые вида $p = c \cdot 2^k + 1$: в $\GF{p}$ есть
    главные корни всех степеней $2^j$, $j \le k$. Примеры: $17 = 2^4 + 1$, $97 = 3 \cdot 2^5 + 1$,
    $257 = 2^8 + 1$, $65537 = 2^{16} + 1$, $998244353 = 119 \cdot 2^{23} + 1$,
    $3221225473 = 3 \cdot 2^{30} + 1$ (первообразные корни $3, 5, 3, 3, 3, 5$). Идея БПФ над
    конечным полем принадлежит Дж.~Полларду~\cite{pollard-fft}; такое преобразование называют
    \emph{теоретико-числовым} (NTT), подробно о нём "--- в лекции~5.
  \item \textbf{Арифметика с плавающей точкой в $\CC$.} Корни $e^{2\pi i k/L}$ и все промежуточные
    значения вычисляются приближённо; результат округляется до ближайшего целого. Округление
    даёт верный ответ, если погрешность каждого $z_k$ меньше $\frac12$.
\end{itemize}

\begin{example}
Для $572 \cdot 38$ (пример~\ref{l4:fft-ex}) годится $p = 257$: $\min(3, 2) \cdot 81 = 162 < 257$
и $4 \mid 256$. Над $\CC$ в двойной точности вычисление из примера~\ref{l4:dft-c} даёт
$z_1 \approx 62 + 1{,}4 \cdot 10^{-15} i$ "--- после округления точно $62$.
\end{example}

\taskexample{fftmul-ex}

\subsection{Точность вычислений с плавающей точкой}

Число двойной точности (IEEE~754) имеет $53$-битную мантиссу, т.\,е. относительную погрешность
$\varepsilon = 2^{-53} \approx 1{,}1 \cdot 10^{-16}$. Обсудим, когда округление заведомо верно
(неформально; точные оценки см. в~\cite[разд.~9.5]{crandall-pomerance},
\cite[разд.~3.3]{brent-zimmermann}).

\begin{itemize}
  \item Сами коэффициенты достигают $\min(M, N)(B - 1)^2 \approx L B^2$, т.\,е.
    $\log_2 L + 2\log_2 B$ бит, и все эти биты должны быть верными: необходимо
    $\log_2 L + 2\log_2 B < 53$.
  \item Каждый уровень БПФ вносит относительную ошибку порядка $\varepsilon$ (плюс ошибку
    вычисленных корней $\omega^k$), уровней $\log_2 L$. В худшем случае абсолютная погрешность
    коэффициентов оценивается величиной порядка $C \cdot \varepsilon \log_2 L \cdot L B^2$ с небольшой
    константой $C$; требование «меньше $\frac12$» даёт запас ещё в несколько бит:
    грубо, $\log_2 L + 2\log_2 B + \log_2\log_2 L + \text{const} \le 53$.
  \item \emph{Сбалансированные цифры} (лекция~2), $\abs{x_i} \le B/2$, уменьшают оценку
    $\abs{z_k} \le \min(M, N) B^2/4$ вчетверо. Важнее то, что слагаемые $x_i y_j$ разных знаков
    в основном взаимно уничтожаются: типичный размер $z_k$ "--- порядка $\sqrt{L}\, B^2$,
    а не $L B^2$, и типичные ошибки ведут себя как случайное блуждание. На практике это позволяет
    брать заметно большие $B$ и $L$, чем даёт гарантированная оценка худшего случая.
\end{itemize}

Например, при умножении двух случайных $2048$-значных десятичных чисел ($L = 4096$, $B = 10$)
максимальное отклонение вычисленных $z_k$ от целых в двойной точности оказывается порядка
$10^{-11}$ "--- огромный запас. Но с ростом $L$ точность $53$ бит неизбежно кончается: чтобы
умножать сколь угодно длинные числа, нужно либо увеличивать точность (и тогда операции над
комплексными числами перестают быть «элементарными»), либо переходить к точной арифметике
и китайской теореме об остатках по нескольким простым $p$ (лекция~5).

\checkyourself{Почему при сбалансированных цифрах и вычислениях в $\GF{p}$ нужно
$p > 2 \max_k \abs{z_k}$, а не $p > \max_k \abs{z_k}$?}

\subsection{Сложность}

\begin{theorem}\label{l4:fftmul}
Пусть числа $a$, $b$ имеют не более $N$ цифр по основанию $B$, $L = 2^r$ "--- наименьшая
степень двойки, не меньшая $2N - 1$, а $p$ "--- простое, $L \mid p - 1$,
$p > N (B - 1)^2$. Тогда алгоритм~\ref{l4:alg-fftmul} в $\GF{p}$ вычисляет $ab$, выполняя
$\bigO(N \log N)$ операций в $\GF{p}$ и $\bigO(N)$ операций над цифрами (переносы).
Точнее: $\frac32 L \log_2 L + \bigO(L)$ умножений и $3 L \log_2 L$ сложений в $\GF{p}$.
\end{theorem}

\begin{proof}
Корректность: по следствию~\ref{l4:linconv} алгоритм вычисляет вычеты $z_k \bmod p$
коэффициентов линейной свёртки; так как $0 \le z_k \le N(B-1)^2 < p$, вычеты совпадают с самими
$z_k$, и по утверждению~\ref{l4:intmul} переносы дают $ab$. ($L$ обратимо в $\GF{p}$, так как
$p$ нечётно.)

Сложность: три БПФ длины $L$ "--- по теореме~\ref{l4:fft} $3 \cdot \frac L2 \log_2 L$ умножений и
$3 L \log_2 L$ сложений; поточечное произведение и умножение на $L^{-1}$ "--- $2L$ умножений;
таблицы степеней $\omega$ и $\omega^{-1}$ "--- $\bigO(L)$; переносы "--- $\bigO(L)$ по
утверждению~\ref{l4:intmul}. Так как $2N - 1 \le L < 4N$, всё это $\bigO(N \log N)$.
\end{proof}

\begin{remark}
Теорема считает операции в $\GF{p}$, а не битовые операции. При фиксированных $p$ и $B$
длина $N$ ограничена условиями $L \mid p - 1$ и $N(B-1)^2 < p$; для растущих $N$ приходится
брать $p$ (или точность в $\CC$) из $\Theta(\log N)$ бит, и битовая сложность становится
$\bigO(N \log N \cdot \mathsf M(\log N))$, где $\mathsf M(k)$ "--- сложность умножения $k$-битных
чисел. Применяя ту же идею рекурсивно, получаем алгоритмы, описанные ниже. Возведение
в квадрат требует лишь двух БПФ вместо трёх.
\end{remark}

\subsection{Быстрее, чем \texorpdfstring{$N \log N \log\log N$}{N log N log log N}}

\begin{itemize}
  \item А.~Шёнхаге и Ф.~Штрассен (1971, \cite{schonhage-strassen}) избавились от проблемы
    точности, выполняя БПФ в кольце $\Zn{(2^K + 1)}$: там $2$ "--- главный корень степени $2K$, и
    умножение на степени $\omega$ "--- это сдвиги. Рекурсивное применение даёт
    $\bigO(n \log n \log\log n)$ битовых операций для $n$-битных чисел
    (\cite[разд.~9.5.6]{crandall-pomerance}, \cite[гл.~8]{von-zur-gathen}). Именно этот алгоритм
    (с улучшениями) используется в GMP для очень больших чисел.
  \item М.~Фюрер (2007) улучшил оценку до $n \log n \cdot 2^{\bigO(\log^* n)}$.
  \item Д.~Харви и Й.~ван~дер~Хувен (2019, \cite{harvey-van-der-hoeven}) построили алгоритм
    сложности $\bigO(n \log n)$ "--- по-видимому, оптимальной; практического значения он
    не имеет из-за огромных констант.
\end{itemize}

\checkyourself{Оцените, во сколько раз БПФ-умножение (теорема~\ref{l4:fftmul}) быстрее
алгоритма Карацубы для $N = 2^{20}$ цифр, считая константы равными единице.}

\checkyourselfsection

\begin{furtherreading}
  \readitem[разд.~9.5.2--9.5.8]{crandall-pomerance}{БПФ, теория свёрток, взвешенные
    и теоретико-числовые преобразования, алгоритм Шёнхаге, сложность умножения}
  \readitem[гл.~30]{clrs}{многочлены и БПФ: интерполяция в комплексных корнях из единицы,
    итеративная реализация}
  \readitem[разд.~4.3.3]{knuth-taocp2}{дискретное преобразование Фурье и алгоритм
    Шёнхаге"--~Штрассена}
  \readitem[гл.~8]{von-zur-gathen}{главные корни из единицы, БПФ и алгоритм
    Шёнхаге"--~Штрассена в общем виде}
  \readitem[разд.~2.3, 3.3]{brent-zimmermann}{преобразование Фурье над кольцами и оценки
    погрешности БПФ с плавающей точкой}
  \readitem{cooley-tukey}{оригинальная статья об алгоритме БПФ}
  \readitem{pollard-fft}{БПФ над конечным полем}
  \readitem{mit-6046-fft}{видеолекция MIT о БПФ и умножении многочленов}
  \readitem{cp-algorithms}{статья «Fast Fourier transform»: реализация БПФ и NTT}
\end{furtherreading}
\end{lecture}
